% paper/sections/appendix_ccd_h01.tex
%
% §B.4  H-01 診断の歴史的ケーススタディ
%   初版（FVM 共存前提）における FCCD 動機付けの保全．
%   本稿の純 FCCD 構成では FVM が消滅するため以下は歴史的経緯である．
%
% ═══════════════════════════════════════════════════════════════════════════════
\section{H-01 診断の歴史的ケーススタディ}
\label{sec:fccd_motivation_h01}
% ═══════════════════════════════════════════════════════════════════════════════

\begin{remark}[本節の位置付け]
本稿の純 FCCD 構成では FVM 面平均演算子 $\mathcal{G}^{\text{adj}}$ は
設計から除外され，{\FCCD} は「面ジェット共通プリミティブ」
（§~\ref{sec:fccd_def}，4 役割 (R1)-(R4)）として再解釈される．
したがって以下の H-01 診断は初版（FVM 共存前提）における歴史的な
動機付けであり，本稿本体の構造的正当化とは独立の経緯記述である．
\end{remark}

\paragraph{H-01 仮説の形成．}
FVM 面平均勾配 $\mathcal{G}^{\text{adj}}$ と
節点中心 CCD 微分の\textbf{メトリック不整合}を先行予備解析が指摘し，
毛管ベンチマーク（$\rho_\text{liq}/\rho_\text{gas}=833$）の Exp-1 で
Balanced-Force 残差 $|\BF_{\text{res}}|\approx 884$（step 1）を実測．
Exp-2（$\sigma=0 \Rightarrow T=20$ まで安定）が決定的証拠となり，
hypothesis \textbf{H-01}（「$\mathcal{G}^{\text{adj}}$ と $\sigma\kappa\nabla\psi$ の
メトリック空間不一致が長時間破綻の主因」）が確定した．

\paragraph{初版（FVM 共存前提）の救済策．}
$\mathcal{G}^{\text{adj}}$ が面上で定義される事実を出発点として，
\textbf{圧力勾配も界面力勾配も同一面ロカス上}で {\FCCD} により評価することで，
BF 残差を $\Ord{\Delta x^2} \to \Ord{\Delta x^4}$ まで圧縮する，というのが
初版（FVM 共存前提）の結論であった．

\paragraph{純 FCCD 設計への昇格．}
本稿では FVM そのものが設計から除外される．
代わりに {\FCCD} は HFE 上流点再構築，GFM ghost pressure jet，
分相 PPE 楕円演算子 $D_h^{bf}$，粘性 stress-divergence 界面バンド
という 4 役割 (R1)-(R4) の\textbf{共通面ジェットプリミティブ}として位置付けられる
（詳細は §~\ref{sec:fccd_def}）．
H-01 は 6 モード障害診断（F-1..F-6）のうち F-1/F-3/F-5 を主因と指摘したが，
現設計ではこれら 3 モードは face-jet 契約の段階で構造的に回避される．

\medskip
\noindent\textbf{補足．}
本付録は H-01 診断の要約的保全であり，
論文本体は (R1)-(R4) による構造的正当化のみで {\FCCD} の必要性を閉じる．

% ═══════════════════════════════════════════════════════════════════════════════
